% GENERATED FILE. Edit canonical lessons, then run npm run generate:books.
% canonical-source-hash: fnv1a64:642cae594859fdda
% canonical-lessons: GU-C348-nyayadhis, GU-C348-kes, GU-C348-apil, GU-C348-pheslo, GU-C348-phatkarvu

\chapter{A Case in Court}
\label{ch:gu-a-case-in-court}
\hlchaptermodality{348}

\begin{chapteropening}
\textbf{By the end of this chapter:} \emph{I can follow what happens to a court case.}

Complete the last lesson of chapter 348: I can follow what happens to a court case.
\end{chapteropening}

\section[\texorpdfstring{\gu{ન્યાયાધીશ} (nyāyādhīś)}{ન્યાયાધીશ (nyāyādhīś)}]{\gu{ન્યાયાધીશ} (nyāyādhīś) — a judge}

\label{lesson:GU-C348-nyayadhis}



\begin{quote}
Before the new one: say the Gujarati for in force, then the Gujarati for to defy.
\end{quote}

\subsection*{You'll want to know: \gu{ન્યાયાધીશ}}
\textbf{\gu{ન્યાયાધીશ}} — \emph{nyāyādhīś} — \textquotedblleft{}a judge\textquotedblright{}.

\begin{quote}
\gu{ન્યાયાધીશે} \gu{બંને} \gu{બાજુની} \gu{વાત} \gu{ધ્યાનથી} \gu{સાંભળી}.
\end{quote}

\begin{quote}
\emph{nyāyādhīśe banne bājunī vāt dhyānthī sā̃bhḷī.}
\end{quote}

\begin{quote}
The judge listened carefully to both sides.
\end{quote}

A masculine or feminine noun, by the person meant.

\begin{cousinweb}
\textbf{\gu{ન્યાયાધીશ}} (\emph{nyāyādhīś}) is built from Sanskrit \emph{nyāya}, \textquotedblleft{}rule, justice\textquotedblright{}, and \emph{adhīśa}, \textquotedblleft{}a lord, a master\textquotedblright{} (Monier-Williams): the master of justice.
\end{cousinweb}

\begin{grammarlens}[title={What you've built}]
One more word for giving reasons and opinions.
\end{grammarlens}

\subsection*{Your turn}
\begin{itemize}
\raggedright
  \item \emph{Say it:} \emph{nyāyādhīś}
  \item \emph{Say it:} \emph{nyāyādhīś}, once more
  \item \emph{Say it:} \emph{nyāyādhīśe banne bājunī vāt dhyānthī sā̃bhḷī.} — The judge listened carefully to both sides.
  \item \emph{Recall:} say the Gujarati for in force, then the Gujarati for to defy, then say \emph{nyāyādhīś} again
\end{itemize}

\begin{tcolorbox}[breakable,colback=teal!4,colframe=teal!35!black,title={Before you move on}]
What does \gu{ન્યાયાધીશ} mean? (\textquotedblleft{}A judge\textquotedblright{}.) Say it once more.
\end{tcolorbox}
\section[\texorpdfstring{\gu{કેસ} (kes)}{કેસ (kes)}]{\gu{કેસ} (kes) — a court case}

\label{lesson:GU-C348-kes}



\begin{quote}
Before the new one: say the Gujarati for to defy, then the Gujarati for a judge.
\end{quote}

\subsection*{You'll want to know: \gu{કેસ}}
\textbf{\gu{કેસ}} — \emph{kes} — \textquotedblleft{}a court case\textquotedblright{}.

\begin{quote}
\gu{આ} \gu{કેસ} \gu{ત્રણ} \gu{વરસથી} \gu{ચાલે} \gu{છે}.
\end{quote}

\begin{quote}
\emph{ā kes traṇ varasthī cāle che.}
\end{quote}

\begin{quote}
This case has been going on for three years.
\end{quote}

A masculine noun.

\begin{cousinweb}
\textbf{\gu{કેસ}} (\emph{kes}) is borrowed from English \emph{case}, in its legal sense.
\end{cousinweb}

\begin{grammarlens}[title={What you've built}]
One more word for giving reasons and opinions.
\end{grammarlens}

\subsection*{Your turn}
\begin{itemize}
\raggedright
  \item \emph{Say it:} \emph{kes}
  \item \emph{Say it:} \emph{kes}, once more
  \item \emph{Say it:} \emph{ā kes traṇ varasthī cāle che.} — This case has been going on for three years.
  \item \emph{Recall:} say the Gujarati for to defy, then the Gujarati for a judge, then say \emph{kes} again
\end{itemize}

\begin{tcolorbox}[breakable,colback=teal!4,colframe=teal!35!black,title={Before you move on}]
What does \gu{કેસ} mean? (\textquotedblleft{}A court case\textquotedblright{}.) Say it once more.
\end{tcolorbox}
\section[\texorpdfstring{\gu{અપીલ} (apīl)}{અપીલ (apīl)}]{\gu{અપીલ} (apīl) — an appeal (to a higher court)}

\label{lesson:GU-C348-apil}



\begin{quote}
Before the new one: say the Gujarati for a judge, then the Gujarati for a court case.
\end{quote}

\subsection*{You'll want to know: \gu{અપીલ}}
\textbf{\gu{અપીલ}} — \emph{apīl} — \textquotedblleft{}an appeal (to a higher court)\textquotedblright{}.

\begin{quote}
\gu{અરુણે} \gu{આ} \gu{કેસમાં} \gu{અપીલ} \gu{કરી}.
\end{quote}

\begin{quote}
\emph{Aruṇe ā kesmā̃ apīl karī.}
\end{quote}

\begin{quote}
Arun appealed in this case.
\end{quote}

A feminine noun.

\begin{cousinweb}
\textbf{\gu{અપીલ}} (\emph{apīl}) is borrowed from English \emph{appeal}.
\end{cousinweb}

\begin{grammarlens}[title={What you've built}]
One more word for giving reasons and opinions.
\end{grammarlens}

\subsection*{Your turn}
\begin{itemize}
\raggedright
  \item \emph{Say it:} \emph{apīl}
  \item \emph{Say it:} \emph{apīl}, once more
  \item \emph{Say it:} \emph{Aruṇe ā kesmā̃ apīl karī.} — Arun appealed in this case.
  \item \emph{Recall:} say the Gujarati for a judge, then the Gujarati for a court case, then say \emph{apīl} again
\end{itemize}

\begin{tcolorbox}[breakable,colback=teal!4,colframe=teal!35!black,title={Before you move on}]
What does \gu{અપીલ} mean? (\textquotedblleft{}An appeal (to a higher court)\textquotedblright{}.) Say it once more.
\end{tcolorbox}
\section[\texorpdfstring{\gu{ફેંસલો} (phẽslo)}{ફેંસલો (phẽslo)}]{\gu{ફેંસલો} (phẽslo) — a ruling}

\label{lesson:GU-C348-pheslo}



\begin{quote}
Before the new one: say the Gujarati for a court case, then the Gujarati for an appeal.
\end{quote}

\subsection*{You'll want to know: \gu{ફેંસલો}}
\textbf{\gu{ફેંસલો}} — \emph{phẽslo} — \textquotedblleft{}a ruling\textquotedblright{}.

\begin{quote}
\gu{ન્યાયાધીશે} \gu{આજે} \gu{કેસનો} \gu{ફેંસલો} \gu{આપ્યો}.
\end{quote}

\begin{quote}
\emph{nyāyādhīśe āje kesno phẽslo āpyo.}
\end{quote}

\begin{quote}
The judge gave the ruling in the case today.
\end{quote}

A masculine noun.

\begin{cousinweb}
\textbf{\gu{ફેંસલો}} (\emph{phẽslo}) is usually traced to Arabic \emph{faiṣala}, \textquotedblleft{}a decision\textquotedblright{}, through Persian.
\end{cousinweb}

\begin{grammarlens}[title={What you've built}]
One more word for giving reasons and opinions.
\end{grammarlens}

\subsection*{Your turn}
\begin{itemize}
\raggedright
  \item \emph{Say it:} \emph{phẽslo}
  \item \emph{Say it:} \emph{phẽslo}, once more
  \item \emph{Say it:} \emph{nyāyādhīśe āje kesno phẽslo āpyo.} — The judge gave the ruling in the case today.
  \item \emph{Recall:} say the Gujarati for a court case, then the Gujarati for an appeal, then say \emph{phẽslo} again
\end{itemize}

\begin{tcolorbox}[breakable,colback=teal!4,colframe=teal!35!black,title={Before you move on}]
What does \gu{ફેંસલો} mean? (\textquotedblleft{}A ruling\textquotedblright{}.) Say it once more.
\end{tcolorbox}
\section[\texorpdfstring{\gu{ફટકારવું} (phaṭkārvũ)}{ફટકારવું (phaṭkārvũ)}]{\gu{ફટકારવું} (phaṭkārvũ) — to hit hard; to impose (a fine or a sentence)}

\label{lesson:GU-C348-phatkarvu}



\begin{quote}
Before the new one: say the Gujarati for an appeal, then the Gujarati for a ruling.
\end{quote}

\subsection*{You'll want to know: \gu{ફટકારવું}}
\textbf{\gu{ફટકારવું}} — \emph{phaṭkārvũ} — \textquotedblleft{}to hit hard; to impose (a fine or a sentence)\textquotedblright{}.

\begin{quote}
\gu{છોકરાએ} \gu{બેટથી} \gu{દડો} \gu{મેદાનની} \gu{બહાર} \gu{ફટકાર્યો}.
\end{quote}

\begin{quote}
\emph{chokrāe beṭthī daḍo medānnī bahār phaṭkāryo.}
\end{quote}

\begin{quote}
The boy whacked the ball out of the ground with the bat.
\end{quote}

Past, with \gu{તેણે} (he or she …): \gu{તેણે} \gu{દડો} \gu{દૂર} \gu{ફટકાર્યો}, \emph{teṇe daḍo dūr phaṭkāryo}.

\begin{cousinweb}
\textbf{\gu{ફટકારવું}} (\emph{phaṭkārvũ}) CDIAL lists the Gujarati verb, \textquotedblleft{}to slash\textquotedblright{}, under \emph{phaṭ}, \textquotedblleft{}a sudden movement\textquotedblright{}, a sound-word like \textquotedblleft{}crack!\textquotedblright{} (CDIAL 9038); a court that \textquotedblleft{}hits\textquotedblright{} someone with a fine uses the same verb.
\end{cousinweb}

\begin{grammarlens}[title={What you've built}]
One more word for giving reasons and opinions.
\end{grammarlens}

\subsection*{Your turn}
\begin{itemize}
\raggedright
  \item \emph{Say it:} \emph{phaṭkārvũ}
  \item \emph{Say it:} \emph{phaṭkārvũ}, once more
  \item \emph{Say it:} \emph{chokrāe beṭthī daḍo medānnī bahār phaṭkāryo.} — The boy whacked the ball out of the ground with the bat.
  \item \emph{Recall:} say the Gujarati for an appeal, then the Gujarati for a ruling, then say \emph{phaṭkārvũ} again
\end{itemize}

\begin{tcolorbox}[breakable,colback=teal!4,colframe=teal!35!black,title={Before you move on}]
What does \gu{ફટકારવું} mean? (\textquotedblleft{}To hit hard; to impose (a fine or a sentence)\textquotedblright{}.) Say it once more.
\end{tcolorbox}
